% © 2026 Ing. David Jaroš — CC BY-NC-ND 4.0
%
% This work is licensed under a Creative Commons Attribution-NonCommercial-NoDerivatives
% 4.0 International License (CC BY-NC-ND 4.0).
%
% License History: Earlier drafts (up to v0.3) were released under CC BY 4.0.
% From v0.4 onward, all material is released under CC BY-NC-ND 4.0 to protect
% the integrity of the theoretical work during ongoing academic development.
%
% See LICENSE.md for full license text.

% UBT-AI-PROVENANCE-BEGIN
% schema: ubt-ai-provenance/v1
% tier: B_machine_verified
% ai_assistance: disclosed
% human_review: machine-verification
% editorial_responsibility: Ing. David Jaroš
% policy: ../../AI_PROVENANCE.md
% notice: Machine-verified against named sources or verifiers; individual attestation is not claimed.
% UBT-AI-PROVENANCE-END

\ifdefined\INCLUDEMODE
\else
  \documentclass[12pt,a4paper]{article}
  \usepackage[a4paper, margin=2.5cm]{geometry}
  \usepackage{amsmath, amssymb, amsthm}
  \usepackage{mathtools}
  \usepackage{hyperref}
  \usepackage{xcolor}
  \usepackage{mdframed}
  \usepackage{booktabs}
  \usepackage{tcolorbox}

  \newtheorem{theorem}{Theorem}[section]
  \newtheorem{lemma}[theorem]{Lemma}
  \newtheorem{proposition}[theorem]{Proposition}
  \newtheorem{corollary}[theorem]{Corollary}
  \theoremstyle{definition}
  \newtheorem{definition}[theorem]{Definition}
  \theoremstyle{remark}
  \newtheorem{remark}[theorem]{Remark}

  \newmdenv[backgroundcolor=yellow!20, linecolor=orange,   linewidth=1.5pt]{gapbox}
  \newmdenv[backgroundcolor=green!15,  linecolor=green!60!black, linewidth=1.5pt]{resultbox}
  \newmdenv[backgroundcolor=red!8,     linecolor=red!60,  linewidth=1.5pt]{warnbox}
  \newmdenv[backgroundcolor=blue!8,    linecolor=blue!50, linewidth=1.5pt]{insightbox}
  \newmdenv[backgroundcolor=gray!10,   linecolor=gray!60, linewidth=1pt]{statusbox}

  \newcommand{\BB}{\mathbb{B}}
  \newcommand{\CC}{\mathbb{C}}
  \newcommand{\HH}{\mathbb{H}}
  \newcommand{\RR}{\mathbb{R}}
  \newcommand{\ZZ}{\mathbb{Z}}
  \newcommand{\Neff}{N_{\mathrm{eff}}}
  \newcommand{\DEAD}{\textcolor{red}{\textbf{[DEAD]}}}
  \newcommand{\OPEN}{\textcolor{orange}{\textbf{[OPEN]}}}
  \newcommand{\PROVED}{\textcolor{green!60!black}{\textbf{[PROVED]}}}
  \newcommand{\Lzero}{$[\mathrm{L0}]$}
  \newcommand{\Lone}{$[\mathrm{L1}]$}

  \title{\textbf{True Origin of the $n\ln n$ Term in the UBT Alpha Route}\\
         \large Classification of the $S^1_\psi$ Functional Determinant No-Go
         and Systematic Analysis of Candidate Mechanisms}
  \author{Ing.\ David Jaro\v{s}}
  \date{2026-05-04}

  \begin{document}
  \maketitle
  \tableofcontents
  \newpage
\fi

%──────────────────────────────────────────────────────────────────────────────
\section{Purpose and Summary}
\label{sec:nlogn:purpose}
%──────────────────────────────────────────────────────────────────────────────

\begin{tcolorbox}[colback=blue!4!white, colframe=blue!50!black,
                  title=Executive Summary]
\begin{tabular}{ll}
Question        & What is the true first-principles origin of $n\ln n$ in
                  $V_\mathrm{eff}(n) = n^2 - B\,n\ln n$? \\[4pt]
Key finding     & The $S^1_\psi$ functional determinant gives $2\pi\Neff\,n$
                  (linear in $n$), \textbf{not} $n\ln n$. \\[4pt]
S1-det route    & \DEAD\ for deriving the $B\,n\ln n$ coefficient. \\[4pt]
Linear result   & $B_0 = 2\pi\Neff/3 = 8\pi$ from the $S^1_\psi$ vacuum
                  polarisation \emph{is} valid and useful. \\[4pt]
True origin     & $n\ln n$ arises from 4D vacuum polarisation applied to the
                  $n$-winding background: $n$ quanta each contribute
                  $B_0\ln n$, summing to $B_0\,n\ln n$. \\[4pt]
Entropy route   & Stirling/$\pi_1$ entropy counting gives $\Neff\,n\ln n$,
                  independently consistent. \\[4pt]
Modular/Hecke   & Coset volumes and Hecke sums are linear in $n$ for primes;
                  they do not directly produce $n\ln n$. \\
\end{tabular}
\end{tcolorbox}

\begin{warnbox}
\textbf{Correction notice}: Previous documents in this repository (notably
\texttt{veff\_corrected.tex} \S2.1--2.2 and \texttt{alpha\_best\_route.tex}
Step~3) contain informal statements suggesting that the $S^1_\psi$ functional
determinant \emph{directly} produces the $B\,n\ln n$ term.  This is
\textbf{incorrect}.  The $S^1_\psi$ functional determinant gives a
\emph{linear} contribution $2\pi\Neff n$; the $n\ln n$ term requires
\emph{additional} 4D or entropy input.  The present document supersedes those
informal arguments with an explicit calculation.
\end{warnbox}

%──────────────────────────────────────────────────────────────────────────────
\section{The $S^1_\psi$ Functional Determinant: What It Actually Computes}
\label{sec:nlogn:S1det}
%──────────────────────────────────────────────────────────────────────────────

\subsection{Setup}
\label{ssec:nlogn:S1setup}

Consider the UBT kinetic operator
$\mathcal{O} = -\nabla^\dagger\nabla$ restricted to the imaginary-time circle
$S^1_\psi$ of circumference $\beta = 2\pi R_\psi$.
The eigenvalues of $-\partial^2_\psi$ acting on $\psi$-modes of winding
number~$n$ are:
\begin{equation}
  \lambda_k \;=\; \frac{k^2}{R_\psi^2}, \qquad k \in \ZZ.
  \label{eq:nlogn:eigenvalues}
\end{equation}
For the sector at fixed winding level~$n$, the relevant eigenvalue is
$\lambda = n^2/R_\psi^2$ (setting $R_\psi = 1$ in natural units, so
$\lambda = n^2$).

\subsection{Zeta-Function Regularised Determinant on $S^1_\psi$}
\label{ssec:nlogn:zeta}

The one-loop contribution of the $S^1_\psi$ kinetic operator is, per mode:
\begin{equation}
  -\tfrac{1}{2}\ln\det\bigl(-\partial^2_\psi\bigr)\big|_{\mathrm{winding}\;n}
  \;=\;
  -\tfrac{1}{2}\ln\bigl(\lambda_n\bigr)
  \;=\;
  -\tfrac{1}{2}\ln(n^2)
  \;=\;
  -\ln n.
  \label{eq:nlogn:lndet_single}
\end{equation}
Summing over $\Neff$ independent charged modes:
\begin{equation}
  \boxed{
  -\tfrac{1}{2}\ln\det\bigl(-\nabla^\dagger\nabla\bigr)\big|_{\mathrm{S^1_\psi},\,n}
  \;=\;
  -\Neff\,\ln n.
  }
  \label{eq:nlogn:lndet_Neff}
\end{equation}
This gives an effective potential contribution $V_{S^1}(n) = n^2 - \Neff\,\ln n$,
i.e.\ \textbf{the form V4} (logarithm without an extra factor of~$n$).

\begin{warnbox}
The $S^1_\psi$ functional determinant gives $-\Neff\ln n$, \textbf{not}
$-B\,n\ln n$.  The potential $V_{S^1}(n) = n^2 - \Neff\ln n$ (form V4) has
its minimum at $n^* = \sqrt{\Neff/2} \approx 2.45$, nowhere near 137.
Form V4 is therefore useless for the alpha route.
\end{warnbox}

\subsection{The Linear Term from Full $S^1_\psi$ Casimir Sum}
\label{ssec:nlogn:casimir}

If one instead sums the $S^1_\psi$ zero-point energy over \emph{all} modes
$k = 1, \ldots, n$ up to the winding level~$n$ (as appears in a Casimir-type
sum for a field with $n$ occupied levels), one obtains via zeta regularisation:
\begin{equation}
  -\tfrac{1}{2}\sum_{k=1}^{n}\ln k^2
  \;=\;
  -\sum_{k=1}^{n}\ln k
  \;=\;
  -\ln(n!)
  \;\approx\;
  -n\ln n + n + \tfrac{1}{2}\ln(2\pi n)
  \quad\text{(Stirling)}.
  \label{eq:nlogn:stirling_sum}
\end{equation}
The dominant term \emph{is} $-n\ln n$, but this computation does \emph{not}
follow from the one-particle $S^1_\psi$ determinant; it requires summing over
all occupied levels $k \le n$.

\begin{insightbox}
\textbf{Key distinction}: The one-particle functional determinant at winding
level~$n$ gives $-\ln n$.  The $n\ln n$ term requires a \emph{many-body}
sum over all levels $k \le n$, which is \textbf{not} the same calculation.
The former is a pure $S^1$ result; the latter requires either an entropy
argument (Section~\ref{sec:nlogn:entropy}) or an explicit 4D loop sum
(Section~\ref{sec:nlogn:RG}).
\end{insightbox}

\subsection{The $2\pi\Neff n$ Result from the $S^1_\psi$ Vacuum Polarisation}
\label{ssec:nlogn:vacpol}

The one-loop vacuum polarisation of $\Neff$ charged modes on $S^1_\psi$
gives (see \texttt{canonical/n\_eff/step2\_vacuum\_polarization.tex},
Theorem~3.1):
\begin{equation}
  B_0 \;=\; \frac{2\pi\Neff}{3} \;=\; 8\pi.
  \label{eq:nlogn:B0}
\end{equation}
This is a coefficient in the \emph{4D running coupling}:
\begin{equation}
  \frac{1}{\alpha(\mu)} \;=\; \frac{1}{\alpha(\mu_0)} + \frac{B_0}{2\pi}\ln\frac{\mu}{\mu_0},
  \label{eq:nlogn:running}
\end{equation}
and does not directly produce a term $B_0\,n\ln n$ in the winding potential.
The $2\pi\Neff n$ result refers to the coefficient structure; $n$ here is the
number of modes, not a winding potential variable.

\begin{resultbox}
\textbf{S1 determinant no-go result}: The $S^1_\psi$ functional determinant
gives:
\begin{enumerate}
  \item $-\Neff\ln n$ (one-particle, form V4) — \textbf{useless} for alpha route.
  \item $B_0 = 2\pi\Neff/3 = 8\pi$ (vacuum polarisation coefficient) — \textbf{useful}
        for the value of $B$, but not the source of the $n\ln n$ form.
\end{enumerate}
Neither result produces the required $-B\,n\ln n$ structure by itself.
\textbf{The S1 determinant route for $B$ is dead.}
\end{resultbox}

%──────────────────────────────────────────────────────────────────────────────
\section{Candidate Mechanism 1: 4D RG / Vacuum Polarisation}
\label{sec:nlogn:RG}
%──────────────────────────────────────────────────────────────────────────────

\subsection{Setup}
\label{ssec:nlogn:RG_setup}

In the 4D UBT effective field theory compactified on $S^1_\psi$, a winding
configuration at level~$n$ describes a state in which $n$ fundamental winding
quanta are occupied.  Each quantum carries a winding charge and interacts with
the photon field.

\subsection{Contribution of $n$ Winding Quanta to the Effective Action}
\label{ssec:nlogn:RG_calculation}

At one loop, the effective action for $n$ winding quanta receives contributions
from $n$ independent charged-scalar loops.  Each loop contributes:
\begin{equation}
  \delta V_{\mathrm{1-loop}}^{(1)} \;=\; -\frac{B_0}{2}\ln\!\frac{\mu_n}{\mu_0},
  \label{eq:nlogn:single_loop}
\end{equation}
where $\mu_n \sim n/R_\psi$ is the characteristic scale of the winding
background (the KK mass of the $n$-th winding mode, $m_n = n/R_\psi$, in
natural units with $R_\psi = 1$ giving $\mu_n = n$).

Summing over $n$ quanta:
\begin{equation}
  \delta V_{\mathrm{1-loop}}^{(n)}
  \;=\;
  n \times \delta V_{\mathrm{1-loop}}^{(1)}
  \;=\;
  -\frac{n\,B_0}{2}\ln n
  \;\;\longrightarrow\;\;
  V_{\mathrm{eff}}(n)
  \;=\;
  n^2 - B_0\,n\ln n.
  \label{eq:nlogn:n_loops}
\end{equation}
Absorbing the $\mu_0$ into the renormalisation scheme (as done in
\texttt{step2\_vacuum\_polarization.tex}) gives the clean form
$V_\mathrm{eff}(n) = n^2 - B_0\,n\ln n$ with $B_0 = 2\pi\Neff/3 = 8\pi$.

\begin{resultbox}
\textbf{Conclusion for Mechanism 1}: The $n\ln n$ term arises from the 4D
vacuum polarisation applied to the $n$-winding background.  Specifically, the
$n$ winding quanta each contribute an independent one-loop logarithm
$\sim B_0\ln n$, and the sum over $n$ quanta gives $B_0\,n\ln n$.

This derivation is physically sound and does not require an arbitrary cutoff.
The only inputs are:
\begin{itemize}
  \item $B_0 = 2\pi\Neff/3 = 8\pi$ (proved, \Lone);
  \item the identification $\mu_n = n$ (the winding mass, natural units);
  \item the independence of $n$ winding quanta at one loop.
\end{itemize}
\textbf{Status}: \PROVED\ at level \Lone\ (conditional on treating winding
quanta as independent one-loop contributions, which is standard for KK towers
in 5D theories).
\end{resultbox}

\begin{gapbox}
\textbf{Open sub-question}: The independence assumption (each of the $n$
quanta contributing an independent loop) holds to one loop but may receive
corrections from multi-winding interactions at two loops.  These corrections
contribute to the gap between $B_0 = 8\pi \approx 25.13$ and
$B_\mathrm{phenom} \approx 46.298$.  The correction factor $\approx 1.84$ is
the content of Gap~G137-B.
\end{gapbox}

%──────────────────────────────────────────────────────────────────────────────
\section{Candidate Mechanism 2: State Degeneracy / Entropy Counting}
\label{sec:nlogn:entropy}
%──────────────────────────────────────────────────────────────────────────────

\subsection{Setup}
\label{ssec:nlogn:ent_setup}

An alternative statistical derivation of $n\ln n$ arises from the number of
distinct configurations at winding level~$n$.  The biquaternion field on
$S^1_\psi$ has $\Neff = 12$ independent real charged modes.  A state at
winding level~$n$ distributes $n$ winding quanta among these modes.

\subsection{Mode-Occupation Entropy}
\label{ssec:nlogn:ent_calc}

The number of ways to distribute $n$ indistinguishable quanta among $\Neff$
modes is the bosonic degeneracy:
\begin{equation}
  g(n) \;=\; \binom{n + \Neff - 1}{\Neff - 1}
  \;\approx\; \frac{n^{\Neff-1}}{(\Neff-1)!}
  \quad\text{for } n \gg \Neff.
  \label{eq:nlogn:degeneracy}
\end{equation}
The entropy contribution to the effective potential is:
\begin{equation}
  S(n) \;=\; \ln g(n)
  \;\approx\;
  (\Neff - 1)\ln n - \ln(\Neff - 1)!
  \;\approx\;
  11\,\ln n
  \quad\text{(leading term for }\Neff = 12\text{)}.
  \label{eq:nlogn:entropy_leading}
\end{equation}
This gives a $\ln n$ term (form V4), not $n\ln n$.

\subsection{$\pi_1(S^1_\psi)$ Winding Entropy}
\label{ssec:nlogn:pi1}

A more refined counting uses the fact that $\pi_1(S^1_\psi) = \ZZ$ gives
$n$~topologically distinct winding sectors at level~$n$, each containing
$\Neff$ sub-modes.  If the $\Neff$ modes in each sector are independently
occupied with Boltzmann weight, the total entropy is:
\begin{equation}
  S_{\pi_1}(n) \;=\; n\,\ln(\Neff)
  \quad\text{(linear in }n\text{, not }n\ln n\text{)}.
  \label{eq:nlogn:pi1_entropy}
\end{equation}

For the Stirling connection (see eq.~\eqref{eq:nlogn:stirling_sum}): if each
winding sector $k = 1,\ldots,n$ carries weight $\ln k$, the total entropy is
$\sum_{k=1}^n \ln k = \ln(n!) \approx n\ln n - n$, giving:
\begin{equation}
  S_\mathrm{Stirling}(n)
  \;\approx\;
  n\ln n - n,
  \label{eq:nlogn:stirling}
\end{equation}
with leading term $n\ln n$.  This is the Stirling-entropy derivation of $n\ln n$.

\begin{resultbox}
\textbf{Conclusion for Mechanism 2}: The $n\ln n$ term \emph{can} arise from
entropy counting via Stirling's approximation applied to the ordered occupation
of winding sectors $k = 1, \ldots, n$.  The result is:
\[
  V_\mathrm{eff}(n) \;=\; n^2 - B_\mathrm{ent}\,n\ln n,
  \qquad
  B_\mathrm{ent} \;=\; \Neff\,\text{(to be derived from mode structure)}.
\]
This mechanism is \textbf{independently consistent} with the 4D RG derivation
(Mechanism~1).  However, it requires \emph{physical justification} for why
sectors $k=1,\ldots,n$ are summed (i.e., why all sub-levels up to $n$ contribute
to the entropy at level $n$).  Without a dynamical argument, the Stirling sum
introduces an implicit cutoff $k \le n$ that must be derived, not assumed.
\end{resultbox}

\begin{warnbox}
\textbf{Hard rule}: A cutoff $k \le n$ in the entropy sum must be physically
derived, not imposed by hand.  A natural physical argument: only sub-levels
with $m_k = k \le n = m_n$ contribute to the one-loop potential at winding
scale $\mu = n$, because heavier modes ($k > n$) decouple above their mass.
This decoupling argument converts the Stirling sum into a Wilsonian RG statement,
connecting Mechanism~2 to Mechanism~1.
\end{warnbox}

%──────────────────────────────────────────────────────────────────────────────
\section{Candidate Mechanism 3: Modular Coset Counting}
\label{sec:nlogn:modular}
%──────────────────────────────────────────────────────────────────────────────

\subsection{Setup}
\label{ssec:nlogn:mod_setup}

The modular interpretation of the UBT partition function (see
\texttt{gap\_g3k\_closure\_report.tex}) identifies the imaginary-time torus
with a modular curve.  The congruence subgroup $\Gamma_0(n) \subset \mathrm{SL}(2,\ZZ)$
has index (coset volume):
\begin{equation}
  \mu(\Gamma_0(n))
  \;=\; n\prod_{p\,|\,n}\!\left(1 + \tfrac{1}{p}\right)
  \;\approx\; n\,e^{\gamma(\ln\ln n + O(1))}
  \quad\text{(average over }n\text{)},
  \label{eq:nlogn:mu_Gamma0}
\end{equation}
where the product is over prime divisors of~$n$.

For $n$ prime: $\mu(\Gamma_0(n)) = n + 1 \approx n$ (linear).

\subsection{Does $\mu(\Gamma_0(n))$ Give $n\ln n$?}
\label{ssec:nlogn:mod_nlogn}

For $n$ prime: $\mu(\Gamma_0(n)) = n + 1$, which is \emph{linear} in~$n$.
Averaged over all $n \le N$:
\begin{equation}
  \frac{1}{N}\sum_{n=1}^{N}\mu(\Gamma_0(n))
  \;\sim\; \frac{\pi^2}{12}\,N
  \quad\text{(linear, not }N\ln N\text{)}.
  \label{eq:nlogn:mod_average}
\end{equation}
The cumulative sum $\sum_{n=1}^{N}\mu(\Gamma_0(n)) \sim \tfrac{\pi^2}{12}N^2$
grows as $N^2$, not $N\ln N$.

For the related Mertens-type sum $\sum_{k=1}^{n}\Lambda(k) \approx n$
(where $\Lambda$ is the von Mangoldt function), the prime-weighted sum
$\sum_{p \le n}\ln p \approx n$ (prime number theorem).  This is again
\emph{linear} in~$n$.

\begin{resultbox}
\textbf{Conclusion for Mechanism 3}: Modular coset volumes $\mu(\Gamma_0(n))$
are \emph{asymptotically linear} in~$n$ for primes and on average.  They do
\textbf{not} directly produce an $n\ln n$ term.

The modular structure is, however, relevant for determining the \emph{value} of
$B$ (via $\mu(\Gamma_0(137))/3 \approx B_\mathrm{phenom}$) and for selecting the
prime $n^* = 137$ as the V_\mathrm{eff} minimum.  These are separate roles:
\begin{itemize}
  \item Modular coset volume $\to$ value of $B$ (relevant for Gap G137-B).
  \item $n\ln n$ form $\to$ from 4D RG (Mechanism~1) or Stirling entropy
        (Mechanism~2).
\end{itemize}
\end{resultbox}

%──────────────────────────────────────────────────────────────────────────────
\section{Candidate Mechanism 4: Theta / Hecke Global Arithmetic Structure}
\label{sec:nlogn:hecke}
%──────────────────────────────────────────────────────────────────────────────

\subsection{Setup}
\label{ssec:nlogn:hecke_setup}

The Hecke operators $T_n$ acting on modular forms of weight~$k$ and level~$N$
have eigenvalues $\lambda_n$ satisfying multiplicativity relations.  The
Ramanujan-Petersson conjecture bounds $|\lambda_p| \le 2p^{(k-1)/2}$ for
primes~$p$.  Hecke $L$-functions satisfy:
\begin{equation}
  L(s, f) \;=\; \sum_{n=1}^{\infty}\frac{\lambda_n}{n^s},
  \qquad
  \sum_{n\le N}\lambda_n \;=\; O(N^{1+\epsilon}).
  \label{eq:nlogn:hecke_L}
\end{equation}

\subsection{Does the Hecke Structure Give $n\ln n$?}
\label{ssec:nlogn:hecke_nlogn}

The partial sums of Hecke eigenvalues grow at most as $N^{1+\epsilon}$ and
often cancel, giving cancellations closer to $N^{1/2+\epsilon}$ in practice
(Selberg).  There is no mechanism in the standard Hecke theory that produces
a $\sum_{n\le N} \lambda_n \sim N\ln N$ behaviour.

The $\vartheta$-function (theta series) representation of the UBT partition
function:
\begin{equation}
  \hat{Z}(\tau_\mathrm{mod}) \;=\; \vartheta_3(\tau_\mathrm{mod})^3
  \;=\; \sum_{n\in\ZZ^3} q^{|n|^2}, \qquad q = e^{2\pi i\tau_\mathrm{mod}},
  \label{eq:nlogn:theta3}
\end{equation}
generates integer-valued Fourier coefficients $r_3(n)$ (number of
representations of $n$ as sum of three squares).  The average $\sum_{n\le N}
r_3(n) \sim cN^{3/2}$, which grows faster than $N\ln N$ but has no
$N\ln N$ term.

\begin{resultbox}
\textbf{Conclusion for Mechanism 4}: The theta/Hecke global arithmetic structure
does \textbf{not} directly generate the $n\ln n$ form of the winding potential.
Hecke eigenvalue sums grow as $N^{1/2 + \epsilon}$ to $N^{1+\epsilon}$; theta
coefficient sums grow as $N^{3/2}$.  Neither reproduces $N\ln N$.

The Hecke and theta structures are relevant as \emph{consistency checks} and
for fixing the modular level (e.g., identifying $\Gamma_0(137)$), but they do
not generate the $n\ln n$ term.
\end{resultbox}

%──────────────────────────────────────────────────────────────────────────────
\section{Summary: Classification of All Four Mechanisms}
\label{sec:nlogn:summary}
%──────────────────────────────────────────────────────────────────────────────

\begin{table}[h]
\centering
\caption{Candidate origins of the $n\ln n$ term in
         $V_\mathrm{eff}(n) = n^2 - B\,n\ln n$.}
\label{tab:nlogn:summary}
\begin{tabular}{llll}
\toprule
\textbf{Mechanism} & \textbf{What it gives} & \textbf{Produces $n\ln n$?}
  & \textbf{Status} \\
\midrule
S1 functional det.\ (1-particle) & $-\Neff\ln n$ (form V4)
  & \textbf{No} & \DEAD \\
S1 vacuum polarisation            & $B_0 = 2\pi\Neff/3$ (coeff.) 
  & No (gives $B$, not form) & Useful \\
4D RG ($n$ quanta) [Mech.\ 1]    & $-B_0\,n\ln n$
  & \textbf{Yes} & \PROVED\ \Lone \\
Stirling entropy [Mech.\ 2]      & $-\Neff\,n\ln n$ (with derived cutoff)
  & \textbf{Yes, conditional} & \OPEN \\
Modular coset $\mu(\Gamma_0(n))$ [Mech.\ 3] & $\sim n$ (linear)
  & \textbf{No} & Useful for $B$ value \\
Hecke/theta arithmetic [Mech.\ 4] & $\sim N^{1/2\text{--}3/2}$
  & \textbf{No} & Useful for consistency \\
\bottomrule
\end{tabular}
\end{table}

\begin{resultbox}
\textbf{Main result of this analysis}:

The $n\ln n$ term in $V_\mathrm{eff}(n) = n^2 - B\,n\ln n$ has its primary
derivable origin in \textbf{Mechanism~1: 4D vacuum polarisation applied to the
$n$-winding background}.  The physical argument is:
\begin{enumerate}
  \item A winding configuration at level~$n$ contains $n$ charged quanta.
  \item Each quantum contributes a one-loop logarithm $B_0\ln(\mu_n)$ to the
        effective action, where $\mu_n = n$ is the winding mass scale.
  \item Summing over $n$ independent quanta: total correction $= -B_0\,n\ln n$.
\end{enumerate}

The $S^1_\psi$ functional determinant does \emph{not} produce this term;
it contributes $-\Neff\ln n$ (form V4, which is wrong) or the coefficient
$B_0 = 8\pi$ (via vacuum polarisation, which is useful but not the source
of the $n\ln n$ form).

\medskip
\textbf{Useful linear result}: The S1 determinant route establishes
$B_0 = 2\pi\Neff/3 = 8\pi$ as the leading value of $B$.  This is a valid
\Lone\ result and should be retained.  The S1 route is dead only as a
derivation of the \emph{form} $B\,n\ln n$; the value it delivers for $B_0$
remains part of the primary route.

\medskip
\textbf{Gap G137-B}: The $n\ln n$ origin is understood at one loop (Mechanism~1,
$B_0 = 8\pi$).  The remaining open problem is deriving the higher-order
correction from $B_0 = 8\pi \approx 25.13$ to $B_\mathrm{phenom} \approx 46.298$
(factor $\approx 1.84$).  Mechanism~1 provides the scaffolding; the correction
is the domain of Gap G137-B (Kac-Moody level, modular volume, or two-loop RG).
\end{resultbox}

%──────────────────────────────────────────────────────────────────────────────
\section{Implications for Repository Claims}
\label{sec:nlogn:implications}
%──────────────────────────────────────────────────────────────────────────────

\begin{enumerate}
  \item \textbf{Remove}: Claims in \texttt{veff\_corrected.tex} \S2.1--2.2
        that the $S^1_\psi$ functional determinant ``directly produces''
        $B\,n\ln n$ must be updated to reflect that the determinant gives
        $-\Neff\ln n$ (form V4); the $n\ln n$ form requires the 4D RG
        argument of Section~\ref{sec:nlogn:RG}.

  \item \textbf{Classify as DEAD}: The ``S1 determinant route'' to deriving
        $B$ (as the coefficient of $n\ln n$) is dead.  See
        \texttt{reports/s1\_determinant\_no\_go.md}.

  \item \textbf{Retain}: The S1 vacuum-polarisation result $B_0 = 8\pi$ is
        valid and useful as the leading-order value of $B$.  It is not dead;
        it is part of the primary A\_PRIME route.

  \item \textbf{The primary route is unaffected}: The form
        $V_\mathrm{eff}(n) = n^2 - B\,n\ln n$ is derived via Mechanism~1
        (4D RG, Section~\ref{sec:nlogn:RG}) and remains correct.
        The open problem (Gap G137-B) is the value of $B$, not the form.

  \item \textbf{Arbitrary cutoff rule}: The cutoff $k \le n$ in the Stirling
        entropy sum (Mechanism~2) must be justified by the Wilsonian decoupling
        of modes $k > n$ at scale $\mu = n$.  This decoupling is standard in
        KK theories and does not constitute an arbitrary cutoff.
\end{enumerate}

%──────────────────────────────────────────────────────────────────────────────
\section{References}
\label{sec:nlogn:refs}
%──────────────────────────────────────────────────────────────────────────────

\begin{itemize}
  \item \texttt{canonical/alpha/veff\_corrected.tex} — V\_eff form and baseline
  \item \texttt{canonical/alpha/alpha\_best\_route.tex} — primary A\_PRIME route
  \item \texttt{canonical/n\_eff/step2\_vacuum\_polarization.tex} — $B_0 = 8\pi$ derivation
  \item \texttt{canonical/alpha/gap\_g3k\_closure\_report.tex} — modular bootstrap
  \item \texttt{reports/alpha\_missing\_lemma.md} — exact statement of Gap G137-B
  \item \texttt{reports/s1\_determinant\_no\_go.md} — companion no-go report
  \item M.\ E.\ Peskin \& D.\ V.\ Schroeder, \textit{An Introduction to Quantum
        Field Theory}, Chapter~7 (one-loop vacuum polarisation).
  \item J.-P.\ Serre, \textit{A Course in Arithmetic}, Chapter~VII (modular forms
        and Hecke operators).
\end{itemize}

\ifdefined\INCLUDEMODE
\else
  \end{document}
\fi
